\documentclass[10pt,a4paper]{amsart}
\usepackage[margin=19mm]{geometry}
\usepackage{amsmath,amssymb,mathtools}
\usepackage[scr=rsfs,cal=euler]{mathalfa}
\usepackage{xfrac}
\usepackage[unicode,hidelinks,bookmarks=false]{hyperref}
\newcommand{\ie}{i.e.\ }
\newcommand{\cf}{cf.\ }
\newcommand{\resp}{respectively\ }
\newcommand{\Subsection}{\subsection}
\providecommand{\nobreakdash}{\nobreak\hbox{-}}
\setlength{\emergencystretch}{2em}
\multlinegap0pt
\usepackage{amssymb}
\usepackage{mathtools}
\usepackage{bm}
\usepackage{xcolor}
\usepackage{csquotes}
\newtheorem{lemma}{Lemma}
\newcommand{\R}{\mathbb{R}}
\newcommand{\Z}{\mathbb{Z}}
\newcommand{\aff}{\operatorname{aff}}
\newcommand{\vct}[1]{\boldsymbol{#1}}
\title{Counting Lattice Points in Minkowski Sums of Cross Polytopes}
\author{Ziyi Dai, Qilin Hou, Zhiyuan Liu, Warut Thawinrak, Hongyu Wang}
\date{Source: arXiv:2608.16037v2 (CC BY 4.0)}
\begin{document}
\maketitle

\noindent\emph{Source and licence.} Counting Lattice Points in Minkowski Sums of Cross Polytopes. Version arXiv:2608.16037v2 (CC BY 4.0), distributed by arXiv under the Creative Commons Attribution 4.0 International licence, http://creativecommons.org/licenses/by/4.0/, as recorded in the arXiv licence metadata for this exact version and shown on the abstract page https://arxiv.org/abs/2608.16037v2. Full licence notice: \texttt{licenses/arxiv-2608.16037-CC-BY-4.0.txt}.\par\smallskip

\noindent\textit{Editorial scope.} Counted unit: the article's lemma lem:slicing-root-polytope (source0.tex lines 586--598). The article's complete original proof of it is delivered with it (source0.tex lines 599--649, one \texttt{proof} environment of 51 lines). No mathematics is written, repaired or re-derived: the statement and the proof are the article's own lines, in the article's own notation, and no retained line is substituted or re-flowed.  Retained uncounted as the source-local context the proof needs: lines 532--539.  Nothing was omitted from any retained span, so the package carries no omission record.  No retained line contains a citation, so the wrapper carries no bibliography and no bibliography entry.  No retained line contains a figure environment, an image-inclusion command or drawing code, so no image is delivered and the package declares no asset dependency.  Editorial formatting only, in addition: an AMS \texttt{amsart} preamble that restates the article's own declarations and macros used by the retained lines, namely the environments \texttt{lemma} on separate counters, together with the standard packages those lines need.  The retained environments print in this document as Lemma 1 in order of appearance, whereas the article prints the same items with the numbering of its own full document.  The complete native source of the article as distributed in the arXiv e-print for this exact version is bundled unchanged as \texttt{sources/207802/source0.tex}.
\begin{equation}
\vct{p}
=\sum_{(\ell_i,r_j)\in\widehat E}\lambda_{ij}(\vct{e}_i+\vct{f}_j),
\text{ for some }
\lambda_{ij}\ge0 \text{ such that }
\sum_{(\ell_i,r_j)\in\widehat E}\lambda_{ij}=k.
\label{eq:flow}
\end{equation}

\begin{lemma}\label{lem:slicing-root-polytope}
Fix $\vct{\alpha}=(\alpha_0,\ldots,\alpha_m)\in\Z_{\ge0}^{m+1}$ and set $
k:=\alpha_0+\cdots+\alpha_m.$ Let 
\[H_{\vct{\alpha}} :=\{(x_0,\dots, x_m,y_0,\dots, y_n) \in \R^{m+1}\oplus \R^{n+1}\mid (x_0,\dots,x_m) = \vct{\alpha}\}\]
be an affine subspace in $\R^{m+1}\oplus \R^{n+1}.$
Then, the polytope $P_{\widehat G,\vct{\alpha}}:=kR_{\widehat G}\cap H_{\vct{\alpha}}$ is integrally equivalent to $P_{G^+}(\alpha_0,\dots, \alpha_m)$, and consequently
\begin{equation}
|k\mathcal R_{\widehat G}\cap\Z^{m+n+2}|
=\sum_{\substack{\alpha_i\ge0\\\alpha_0+\cdots+\alpha_m=k}}
|P_{G^+}(\alpha_0,\ldots,\alpha_m)\cap\Z^n|.
\label{eq:slicing}
\end{equation}
\end{lemma}

\begin{proof} Let
\[
\pi:\R^{m+1}\oplus\R^{n+1}\longrightarrow\R^{n}
\]
be the projection given by $\pi(x_0,\dots, x_m,y_0,\dots, y_n) := (y_1\dots, y_n)$. We first show that $\pi$ defines an invertible transformation from $\aff(P_{\widehat G,\vct{\alpha}})$ to $\aff(P_{G^+}(\alpha_0,\dots,\alpha_m))$ and preserves the lattice points between $P_{\widehat G,\vct{\alpha}}$ and $P_{G^+}(\alpha_0,\dots,\alpha_m)$. 

Note that every point $\vct{p}$ in $P_{\widehat G,\vct{\alpha}}$
 has the form $\vct{p}=(\alpha_0,\dots,\alpha_m,\beta_0,\dots, \beta_n)$
 for some nonnegative real numbers $\beta_0,\dots, \beta_n.$ By representing $\vct{p}$ as in \eqref{eq:flow}, we define for each $i\in\{0,1,\ldots,m\}$
%its contribution to the genuine right coordinates by
\[
\vct{w}^{(i)}:=(\lambda_{i1},\ldots,\lambda_{in})\in\R_{\ge0}^n.
\]
%For $i\ge1$, this vector is supported on $I_i$, while for $i=0$ it may be supported anywhere in $[n]$. 
Since the $i^{\mathrm{th}}$-row-sum is $\alpha_i$, we have
\[
\sum_{j=1}^n w_j^{(i)}\le\alpha_i.
\]
Therefore,
$\vct{w}^{(0)}\in\alpha_0\Delta_{[n]}^0$ and $\vct{w}^{(i)}\in\alpha_i\Delta_{I_i}^0$ for all $i\in [m]$. Since
\[
\vct{\beta}:=(\beta_1,\ldots,\beta_n)
=\sum_{i=0}^m\vct{w}^{(i)},
\]
it follows that
\[
\vct{\beta}
\in\alpha_0\Delta_{[n]}^0+\sum_{i=1}^m\alpha_i\Delta_{I_i}^0
=P_{G^+}(\alpha_0,\ldots,\alpha_m).
\]
This shows that the projection $\pi$ satisfies $\pi(\vct{p}) = \vct{\beta} \in P_{G^+}(\alpha_0,\dots, \alpha_m)$ for all $\vct{p} \in P_{\widehat G,\vct{\alpha}}$, and hence defines a map from $P_{\widehat G,\vct{\alpha}}$ to $P_{G^+}(\alpha_0,\dots, \alpha_m).$ 

Conversely, suppose
$\vct{\beta}\in P_{G^+}(\alpha_0,\ldots,\alpha_m).$ By the definition of a Minkowski sum, we may choose $\vct{w}^{(0)}\in\alpha_0\Delta_{[n]}^0$ and $
\vct{w}^{(i)}\in\alpha_i\Delta_{I_i}^0$
such that $\vct{\beta}=\sum^m_{i=0}\vct{w}^{(i)}$. Set
\[
\lambda_{ij}:=w_j^{(i)} \text{ for }j\ge1
\quad \text{and } \quad
\lambda_{i0}:=\alpha_i-\sum_{j=1}^n w_j^{(i)}\ge0.
\]

Hence, $\lambda_{ij}$ can be regarded as a nonnegative weight of the edge $(i,j)$ of $\widehat G$. All these weights of the edges of $\widehat G$ have their $i^{\mathrm{th}}$-row-sums equal to $\alpha_i$ for all $i \in [m]$. By letting
\[
\beta_0=k-\sum_{j=1}^n\beta_j,
\]
one sees that $\vct{p}:=(\vct{\alpha},\beta_0,\vct{\beta})$ is a point in $P_{\widehat G, \vct{\alpha}}$ satisfying $\pi(\vct{p}) = \vct{\beta}.$  Clearly, once $\vct{\alpha}$ and $\vct{\beta}$ are fixed, such a point $\vct{p}$ is uniquely determined. Moreover, $\vct{\beta}$ is a lattice point if and only if $\vct{p}$ is a lattice point. This shows that $\pi$ defines an invertible map from  $P_{\widehat G,\vct{\alpha}}$ to $P_{G^+}(\alpha_0,\dots, \alpha_m)$ and preserves their lattice points. The linearity of $\pi$ then implies that it actually defines an invertible affine transformation from $\aff(P_{\widehat G,\vct{\alpha}})$ to $\aff(P_{G^+}(\alpha_0,\dots,\alpha_m))$. Therefore, the two polytopes $P_{\widehat G,\vct{\alpha}}$ and $P_{G^+}(\alpha_0,\dots,\alpha_m)$ are integrally equivalent.


Finally, every lattice point of $k\mathcal R_{\widehat G}$ has a unique integral left marginal $\vct{\alpha}\in\Z_{\ge0}^{m+1}$ with $|\vct{\alpha}|=k$. Thus, summing over their cardinalities yields \eqref{eq:slicing}.
\end{proof}

\end{document}
